\documentclass{book}
\usepackage{graphicx} % Required for inserting images
\usepackage{url, hyperref}
\usepackage[a4paper]{geometry}
\usepackage{verbatim}
\usepackage[normalem]{ulem}
\usepackage{fancyhdr}
\pagestyle{fancy}
\lhead{easyPython}
\chead{Chapter-1}
\rhead{\thepage}
\lfoot{If error anywhere, contact:engilazarus10@gmail.com}

\title{easyPythoN}
\author{Writter: Lazarus Tuwary Tripura}
\date{written from 6/07/2026 }

\begin{document}

\maketitle

\section*{Introduction}
\addcontentsline{toc}{section}{Introduction}
Learning programming feels easy when you only follow simple tutorials like print("Hello, world!") but not easy cause need to pay a price. Classes books and courses teach concepts step by step, but many struggle when they try to create their own programs from zero knowledge.The main problem is this: Book teach syntax and ideas, but not always how to think like a programmer. When students open a blank editor to build something alone, they often do not know where to begin. Project-based programming books are usually far better for developing real engineering ability than only syntax/tutorial learning.\\
\\ Beginners often struggle building large software projects. Small programs are easier to understand and help students learn programming step by step.\\
This book uses simple design principles:
\begin{itemize}
    \item Small Programs
    \item Text-Based
    \item Single File programs
    \item Many Project programs
    \item Friendly Code to Advanced\\\\
The book helps you understand how real programs are built by practicing with full projects. Let us have fun by acquiring knowledge and practice.
\end{itemize}
\textbf{Recommended learning steps:}
\begin{enumerate}
    \item Type the program first and then understand what it does
    \item Don't try to copy-paste and fix if error comes and ask AI, why it become
    \item Follow the steps and learn the basics
    \item Modify the parts and observe the changes
    \item If you have time then try to alternate way has or not
\end{enumerate}
Do you use windows or Linux? But both can download: \href{https://code.visualstudio.com/}{VS Code}.This is an editor so Python itself must be download: \href{https://www.python.org/downloads/}{Python}
Again some pain you have to download extension for VS code: \href{https://marketplace.visualstudio.com/items?itemName=ms-python.python&utm_source}{Extension}
But if you use Ubuntu on linux, Open Gnome terminal and type or copy: \\
sudo apt update\\
sudo apt install python3\\
python3\\
So technically:
\begin{itemize}
    \item You do not need the VS Code
    \item You do not need graphical software
    \item Only the terminal is enough
\end{itemize}
\section*{Installation Python Modules}
\addcontentsline{toc}{section}{Installation Python Modules}
However, some projects need extra modules. They save huge development time.\\
Let's know about modules-

\begin{enumerate}
    \item pyperclip(Python itself does not directly provide easy clipboard control.)\\Example uses: Password Manager, Automation Scripts, Automatically Copy Generated
    \item bext(Normal Python terminal output is plain text only. Python terminal output into colored text, cursor-controlled layouts, and terminal graphics for better readability and user interaction for use like ex:terminal games, dashboards, animations system monitors)
    \item playsound(Python cannot easily play MP3/WAV files using only basic built-in tools. Ex: alarms, AI assistants, notifications, voice systems)
    \item pyttsx3(Python has no built-in speech engine so why does it need to program compile. Ex:voice assistants, accessibility systems, robotics)
\end{enumerate} 
Without libraries, programmers would need to build complex systems themselves, such as audio playback engines, speech synthesis systems, and direct operating system interfaces. This would require enormous time and deep low-level programming knowledge. Libraries solve this problem by providing ready-made tools and functions, allowing developers to focus on solving real problems, building useful products, and creating applications much faster and more efficiently.\\
Copy and Paste
\begin{center}
    "python3 -m pip install bigbookpython"
\end{center}
For Visual Studio Code or Python shell-
\begin{center}
    "import os, sys"\\
    "os.system(sys.executable + ' -m pip install --user bigbookpython')"
\end{center}
To test installation:
\begin{center}
    "import pyperclip"\\
    "import bext"
\end{center}
In Linux or terminal: 
\begin{center}
    "cd ~/python"\\  
    "python3 -m venv .venv"\\
    "pip install bigbookpythonCopied!"   
\end{center}
Hey! There are some difficulties, but you know that you have to pay a price, this is the reality. But don't worry, programming improves through practice, not just reading but real fun having so draw near. It is boring to read; let us solve the real problem and observe.\\
\uline{\textbf{Project 1}}: Clipboard Automation and Smart Copy-Paste Management System using Python and Pyperclip. (Understanding:the standard clipboard into a powerful tool that can monitor, format, filter, and store data automatically.)


\end{document}
